\begin{lemma}[Linearity of the weak gradient, and the constants]
Let $n\ge1$ and $q\ge1$, let $\Omega\subseteq\R^n$ be open, let
$f_1,f_2\colon\R^n\to\R$ and $g_1,g_2\colon\R^n\to\R^n$, and let
$a,b,c\in\R$.  Then:
\begin{enumerate}
\item if $f_1\in W^{1,q}_{\mathrm{loc}}(\Omega)$ with weak gradient $g_1$ and
  $f_2\in W^{1,q}_{\mathrm{loc}}(\Omega)$ with weak gradient $g_2$
  (\noderef{MemW1pLocOn}), then
  $af_1+bf_2\in W^{1,q}_{\mathrm{loc}}(\Omega)$ with weak gradient
  $ag_1+bg_2$;
\item if $f_1\in W^{1,q}(\Omega)$ with weak gradient $g_1$ and
  $f_2\in W^{1,q}(\Omega)$ with weak gradient $g_2$ (\noderef{MemW1pOn}), then
  $af_1+bf_2\in W^{1,q}(\Omega)$ with weak gradient $ag_1+bg_2$;
\item the constant function $x\mapsto c$ lies in
  $W^{1,q}_{\mathrm{loc}}(\Omega)$ with weak gradient $0$; and
\item if $|\Omega|<\infty$, the constant function $x\mapsto c$ lies in
  $W^{1,q}(\Omega)$ with weak gradient $0$.
\end{enumerate}

The four clauses are one node because they are one fact used in one way: every
consumer that forms a difference $u-v$, a translate $u+c$ or a linear
combination of two Sobolev functions needs the corresponding clause, and the
proofs of the four share their only substantial step, the splitting of the
defining integral identity.  Clause (4) needs $|\Omega|<\infty$ and clauses
(1)--(3) do not: a constant $c\ne0$ has $\int_\Omega|c|^q=\infty$ when
$|\Omega|=\infty$, so it is not in $L^q(\Omega)$, whereas on a compact subset of
$\Omega$ --- which is all clause (3) asks about --- the measure is automatically
finite.  Clause (2) needs no finiteness hypothesis either, because the
$L^q(\Omega)$ membership of the combination comes from that of $f_1$ and $f_2$
and not from the measure of $\Omega$.

Nothing here is specific to the operators of the paper: the statement mentions
no $A$, no solution and neither structure condition.  It is the ambient
linearity of distributional differentiation, stated for the project's own
predicates because, by the standing dividing line, facts about
\noderef{HasWeakGradientOn}, \noderef{MemW1pOn} and \noderef{MemW1pLocOn} are
nodes and not Preamble background.
\end{lemma}
\begin{proof}
Throughout, ``test function'' means a function $\varphi\in C^\infty(\R^n)$ with
compact support and $\operatorname{supp}\varphi\subseteq\Omega$, which is
verbatim the test class of \noderef{HasWeakGradientOn}, and $D_i$ denotes the
partial derivative in the direction of the $i$-th standard basis vector.

\emph{Step 1: the four integrals in play are finite.}  Let $\varphi$ be a test
function, let $i\in\{1,\dots,n\}$, and put $K=\operatorname{supp}\varphi$, a
compact subset of $\Omega$.  Suppose first that $f_1,f_2$ satisfy the hypothesis
of clause (1).  By the local clause of \noderef{MemW1pLocOn} applied to the
compact set $K\subseteq\Omega$,
\[
f_1,f_2\in L^q(K),\qquad |g_1|,|g_2|\in L^q(K),
\]
and in particular each component $g_k^i$ is in $L^q(K)$, since
$|g_k^i|\le|g_k|$ pointwise.  As $K$ is compact its measure is finite, so
H\"older on a set of finite measure (Preamble) gives
$L^q(K)\subseteq L^1(K)$ and hence
\[
f_1,f_2,g_1^i,g_2^i\in L^1(K).
\]
The functions $\varphi$ and $D_i\varphi$ are continuous with compact support,
hence bounded on $\R^n$, say by $\Lambda$.  Each of the four functions
$f_kD_i\varphi$ and $g_k^i\varphi$ is measurable, vanishes off $K$ (because
$\varphi$ and $D_i\varphi$ do: $\varphi$ vanishes on the open complement of $K$,
so all its derivatives vanish there too), and is dominated on $K$ by
$\Lambda|f_k|$ respectively $\Lambda|g_k^i|$, which are in $L^1(K)$.  By the
Preamble item on domination, all four are integrable over $\Omega$.

\emph{Step 2: the weak-gradient identity for the combination.}  By
\noderef{HasWeakGradientOn}, which is the first clause of
\noderef{MemW1pLocOn}, we have for $k=1,2$
\[
\int_\Omega f_k\,D_i\varphi\,dx=-\int_\Omega g_k^i\,\varphi\,dx .
\]
Multiply the identity for $k=1$ by $a$ and the one for $k=2$ by $b$ and add.
Every integral occurring is finite by Step 1, so additivity and homogeneity of
the integral apply on both sides and give
\[
\int_\Omega\bigl(af_1+bf_2\bigr)D_i\varphi\,dx
=-\int_\Omega\bigl(ag_1^i+bg_2^i\bigr)\varphi\,dx
=-\int_\Omega\bigl(ag_1+bg_2\bigr)^i\varphi\,dx ,
\]
the last step because the $i$-th component of $ag_1+bg_2$ is $ag_1^i+bg_2^i$.
As $\varphi$ and $i$ were arbitrary, this is exactly the defining identity of
\noderef{HasWeakGradientOn} for the pair $(af_1+bf_2,\ ag_1+bg_2)$ on $\Omega$.

\emph{Step 3: the $L^q$ clauses, and clauses (1) and (2).}  Let
$K'\subseteq\Omega$ be compact.  By the local clause of \noderef{MemW1pLocOn},
$f_1,f_2\in L^q(K')$ and $|g_1|,|g_2|\in L^q(K')$.  By the Preamble item on the
$L^q$ triangle inequality, which also records closure under multiplication by a
real scalar and holds for functions valued in a real Banach space as well as for
real-valued ones, $af_1+bf_2\in L^q(K')$ and $|ag_1+bg_2|\in L^q(K')$.  Together
with Step 2 this is clause (1) (\noderef{MemW1pLocOn}).

For clause (2) the argument is the same with one change.  The hypothesis is now
$f_1,f_2\in L^q(\Omega)$ and $|g_1|,|g_2|\in L^q(\Omega)$ (the second and third
clauses of \noderef{MemW1pOn}).  Step 1 needs $f_k,g_k^i\in L^1(K)$ for the
compact $K=\operatorname{supp}\varphi\subseteq\Omega$; this follows from
$L^q(\Omega)$ membership by restriction to the subset $K$ (Preamble, the
transport facts recorded with H\"older) and then H\"older on $K$, whose measure
is finite.  Step 2 is then unchanged.  The $L^q(\Omega)$ membership of
$af_1+bf_2$ and of $|ag_1+bg_2|$ is again the $L^q$ triangle inequality with
scalars (Preamble), now on $\Omega$ itself.  This is clause (2)
(\noderef{MemW1pOn}).

\emph{Step 4: the constant function has weak gradient $0$.}  Let $\kappa$ denote
the function constantly equal to $c$, let $\varphi$ be a test function and let
$i\in\{1,\dots,n\}$.  Apply the Preamble item on integration by parts in a
direction with $v=e_i$, with $f\equiv1$ the constant function $1$ and with
$g=\varphi$.  Its hypotheses hold: $f$ is differentiable at every point, with
$D_if\equiv0$; $\varphi$ is $C^\infty$, so differentiable at every point; and
the three functions $D_if\cdot\varphi=0$, $f\cdot D_i\varphi=D_i\varphi$ and
$f\varphi=\varphi$ are integrable on $\R^n$, each being continuous with compact
support or identically $0$.  The conclusion is
\[
\int_{\R^n}D_i\varphi\,dx=\int_{\R^n}f\,D_i\varphi\,dx
=-\int_{\R^n}D_if\cdot\varphi\,dx=0 .
\]
Since $D_i\varphi$ vanishes off $\operatorname{supp}\varphi\subseteq\Omega$, the
integral over $\R^n$ equals the integral over $\Omega$: the two integrands
$\mathbf1_{\R^n}D_i\varphi$ and $\mathbf1_\Omega D_i\varphi$ are the same
function.  Hence
\[
\int_\Omega\kappa\,D_i\varphi\,dx=c\int_\Omega D_i\varphi\,dx=0
=-\int_\Omega 0\cdot\varphi\,dx ,
\]
where the first equality is homogeneity of the integral, legitimate because
$D_i\varphi$ is integrable over $\Omega$ (continuous with compact support).
This is the defining identity of \noderef{HasWeakGradientOn} for the pair
$(\kappa,0)$ on $\Omega$.

\emph{Step 5: clauses (3) and (4).}  For clause (3), let $K'\subseteq\Omega$ be
compact; then $|K'|<\infty$, so the constant $\kappa$ lies in $L^q(K')$ and so
does the zero field (Preamble: a constant function lies in $L^q$ of a finite
measure).  With Step 4 this is clause (3) (\noderef{MemW1pLocOn}).  For clause
(4), the hypothesis $|\Omega|<\infty$ gives $\kappa\in L^q(\Omega)$ and
$0\in L^q(\Omega)$ by the same Preamble item, and Step 4 supplies the
weak-gradient identity; this is clause (4) (\noderef{MemW1pOn}).
\end{proof}
